\section{Discussion, limitations and conclusion}
\label{sec:conclusion}

\paragraph{What worked.}
The measurable gains came from three sources:
\begin{itemize}
  \item searching deeper per second: bitboards, incremental distances, move ordering;
  \item removing a systematic bias from the evaluation;
  \item a few targeted rules found by reading replays: the progress tie-break, the rush rule
        and defensive walls.
\end{itemize}
Each of these fixed a concrete, reproducible failure. Learning the evaluation with TD-leaf gave
an interpretable evaluation of equal strength rather than a stronger one. Under this search
regime, adding more static information to the evaluation did not change results (\#35, \#37,
\#38).
Pure self-play has a blind spot: the agent only learns how much a held wall is worth against
itself. Testing against bots with a different style (\#44--45) exposed this, and a one-parameter
correction fixed it.

\paragraph{Limitations and threats to validity.}
\begin{itemize}
  \item \emph{Grading speed is simulated.} The instructor confirmed that grading runs on a free
        Colab CPU and that the one-second limit includes inter-process communication. We assume
        that machine is about $2\times$ slower than our laptop (Intel i5-11400H); this ratio has
        not been measured on Colab itself. Most tests slow down only our agent, which is
        pessimistic for us; the key standard-start tests slow down both agents.
  \item \emph{Noise.} With 40 games the standard error of a score is about 0.08, and several
        decisions rest on differences of 0.05--0.10. Where possible we repeated runs on a second
        seed and pooled 160--240 games.
  \item \emph{Standard start is nearly deterministic.} Two deterministic agents produce only two
        distinct standard-start games, varied only by timing. A fixed losing line could
        therefore exist against an unseen opponent even though random-opening results are good.
  \item \emph{Horizon effect.} The agent still sometimes delays when one wall would hurt it
        badly. Defensive walls cover single-wall traps only, and only at the root.
  \item \emph{Opponent pool.} The pool contains our own earlier versions, so it may favour
        styles similar to ours. External programs (Hard, gorisanson's MCTS) reduce but do not
        remove this bias.
\end{itemize}

\paragraph{Conclusion.}
Our answer to the research question: under a one-second pure-Python budget, the strongest
design we found was a fast, well-ordered alpha-beta search with a symmetric evaluation that is
cheap enough to keep depth high, plus testing against opponents of other styles. The submitted
agent scores 0.87 from the standard start and 0.88 with random openings against 16 opponents,
and 0.775 from the standard start against wrongway.app's Hard AI at simulated grading speed. Its
main known weakness is still the standard start against wall-hoarding bots, where the result
depends on timing.
Depth~8 would win 0.81 against Hard, but two further speed-up attempts (\#42--43) gave no gain,
and a parallel search is not worthwhile on the single physical core of a free Colab CPU. The most
promising future work is therefore to learn the evaluation against a diverse pool of opponents
rather than by self-play alone.
